\documentclass[11pt]{article}
\usepackage[utf8]{inputenc}
\usepackage{amsmath}
\usepackage{amsfonts}
\usepackage{amssymb}
\usepackage{bm}
\usepackage{geometry}
\usepackage{xcolor}
\usepackage{booktabs}
\usepackage{enumitem}

\geometry{margin=0.9in}

\title{Evaluating Recovery of a Dale-Constrained Connectivity Matrix\\
from Synthetic Spike Data:\\
Study Design and Comparison Metrics}
\author{Jan Balewski, NERSC}
\date{\today}

\newcommand{\Atrue}{A_{\mathrm{true}}}
\newcommand{\Ahat}{\hat{A}}
\newcommand{\Etrue}{E_{\mathrm{true}}}
\newcommand{\minW}{\mathrm{minW}}
\newcommand{\neff}{n_{\mathrm{eff}}}
\newcommand{\Trec}{T_{\mathrm{rec}}}

\begin{document}

\maketitle

\begin{abstract}
We define a protocol for quantifying how accurately the shared connectivity
matrix $\Atrue$ of a Dale-law-constrained recurrent Poisson GLM is recovered
by PRISM-EM from non-stationary spike trains, as a function of recording
duration. Ground truth is known by construction. The evaluation is split into
three independent concerns---\emph{support} (which edges exist), \emph{sign}
(excitatory vs.\ inhibitory), and \emph{magnitude} (edge weight)---of which
this document treats support and magnitude in detail. We fix the pruning
threshold $\minW$ and the regularization strength for the present study, apply
$\minW$ to the raw fitted weights \emph{before} any spectral rescaling, and
compare magnitudes only after normalizing both matrices to unit spectral
radius. We motivate the expected $\propto 1/\sqrt{\neff}$ scaling of weight
estimation error with effective sample size and use it as the quantitative
null hypothesis for the duration sweep.
\end{abstract}

%\tableofcontents
%\newpage

% ===============================================================
\section{Background: Data and Fit in Brief}
\label{sec:background}
% ===============================================================

\subsection{Generator}
A network of $N$ neurons ($N_E$ excitatory, $N_I = N - N_E$ inhibitory) is
simulated as a lag-1 Poisson vector autoregression (VAR(1)). A single shared
connectivity matrix $\Atrue \in \mathbb{R}^{N\times N}$ governs all dynamics;
$M$ latent states differ only in a per-state bias (log-baseline) vector $B_m$.
At each time bin $t$ of width $\Delta t$,
\begin{equation}
  \eta_t = \Atrue\, Y_{t-1} + B_{\mathrm{eff},t},\qquad
  \lambda_t = \exp\!\bigl(\min(\eta_t,\eta_{\mathrm{clip}})\bigr)\Delta t,\qquad
  Y_t \sim \mathrm{Poisson}(\lambda_t),
\end{equation}
where $B_{\mathrm{eff},t}=\sum_m c_{t,m}B_m$ mixes the state biases along a
smooth simplex trajectory $c_t$. Column convention: $(\Atrue)_{ij}$ is the
weight \emph{from} presynaptic neuron $j$ \emph{to} postsynaptic neuron $i$.
Dale's law is enforced at the source: every outgoing weight of an excitatory
neuron is positive, every outgoing weight of an inhibitory neuron is negative,
and self-loops are negative. The binary support $\Etrue\in\{0,1\}^{N\times N}$
is generated by per-neuron i.i.d.\ Bernoulli connectivity, and the spectral
radius is fixed at $\rho(\Atrue)=R<1$. Ground-truth arrays
$(\Atrue, \{B_m\}, \Etrue)$ and the $(x,y)$ coordinates of all neurons are
saved.

\subsection{Fitter (PRISM-EM)}
PRISM-EM recovers $\Ahat$ and $\{\hat B_m\}$ by block coordinate descent: an
E-step infers the simplex coefficients $c_{1:T}$ and Viterbi-decodes them to
one-hot state assignments; an M-step updates $\Ahat,\hat{\mathbf B}$ by Adam on
the Poisson negative log-likelihood with off-diagonal $\ell_1$ shrinkage
(strength $\lambda_3$) and a spectral-radius projection to $\rho_{\max}$. The
fitter does \emph{not} impose Dale's law, so recovered rows may contain mixed
signs; this is exploited later only as a sign-consistency diagnostic and is out
of scope here. Because $\Atrue$ is shared across all states, every time bin
contributes to estimating $\Ahat$, and the arbitrary permutation of recovered
state labels does \emph{not} affect $\Ahat$; the connectivity comparison
therefore requires no label alignment.

% ===============================================================
\section{Part I: Study Design and Conventions}
\label{sec:design}
% ===============================================================

\subsection{Operating Regime}
We target $N\in[50,500]$ neurons with a median single-neuron firing rate of
$2$--$3$\,Hz (range $[1,5]$\,Hz), spikes binned at $\Delta t = 10$\,ms, and
recording durations from $10$ to $60$ minutes. The latent dynamics use
$M\in\{2,3\}$ states with mean dwell time $\sim 1$\,s. The central scientific
question is: \emph{how much recovery quality is lost by using a $10$-minute
record instead of a $60$-minute record?}

\subsection{Master-Record Subsampling}
To remove generator variability from the duration comparison, a single
\emph{master} record of $120$ minutes is generated per configuration
$(N,\text{seed})$. Shorter records are obtained as contiguous temporal
prefixes,
\begin{equation}
  Y^{(D)} = Y_{1:T_D},\qquad T_D = \lfloor D/\Delta t\rfloor,
\end{equation}
for durations $D\in\{10,15,20,30,45,60\}$ minutes, with the full $120$-minute
record retained as an asymptotic ceiling. Contiguous prefixes (rather than
random bin subsets) preserve the temporal lag structure and the state-dwell
statistics that the fitter depends on. Each metric is reported relative to its
$120$-minute value, so results read as ``fraction of asymptotic recovery
retained.'' At least $5$ independent seeds are run per $(N,D)$ pair; sampling
variability in this sparse, low-rate regime is large and single-seed results
are not interpretable.

\subsection{Parameters Held Constant (Present Study)}
For this phase the following are fixed across all durations and seeds:
\begin{itemize}[nosep]
  \item Pruning threshold $\minW = 0.02$ (applied to raw $\Ahat$; see below).
  \item $\ell_1$ strength $\lambda_3$ and all other fit hyperparameters
        (learning rates, EM/M-step iteration counts, $\rho_{\max}$, delay
        schedules).
  \item Network structure $(N, N_E, R, \Etrue, \Atrue, \{B_m\})$ within a seed.
\end{itemize}
Holding $\lambda_3$ fixed means we measure \emph{fixed-regularizer} recovery.
Because the optimal $\ell_1$ penalty scales with the noise level
($\propto 1/\sqrt{\neff}$, Section~\ref{sec:neff}), a fixed $\lambda_3$
over-shrinks short records; the measured duration effect therefore includes a
regularization component that we acknowledge explicitly and revisit in later
work.

\subsection{Threshold Then Rescale: Order of Operations}
\label{sec:order}
Support and magnitude use different normalizations of $\Ahat$, and the order
matters.
\begin{enumerate}[nosep]
  \item \textbf{Support} is evaluated on the \emph{raw} fitted matrix. An edge
        $(i,j)$ is declared present iff $|\Ahat_{ij}| > \minW$, $i\neq j$. The
        threshold is absolute and is applied at the fitter's native scale
        ($\rho \approx \rho_{\max}$).
  \item \textbf{Magnitude} is evaluated after rescaling both matrices to unit
        spectral radius,
        \begin{equation}
          \Ahat^{\star} = \Ahat / \rho(\Ahat),\qquad
          \Atrue^{\star} = \Atrue / \rho(\Atrue),
        \end{equation}
        which removes the global scale mismatch between the fitter's
        $\rho_{\max}$ constraint and the generator's $R$. This rescaling is a
        single scalar per matrix; it preserves all signs, the zero pattern, and
        all relative orderings.
\end{enumerate}
The two normalizations are kept in separate pipelines. Applying the absolute
$\minW$ after a $\rho$-rescaling would make its effective cut size depend on
$N$ (since unit-$\rho$ entries scale roughly as the inverse square root of the
number of active edges), introducing an $N$-dependent artifact into the support
results; thresholding the raw weights avoids this. All metrics, both support
and magnitude, exclude the diagonal.

\subsection{Effective Sample Size and the $1/\sqrt{\neff}$ Law}
\label{sec:neff}
The recoverability of an edge $j\!\to\!i$ is governed not by raw duration but
by how often the presynaptic neuron $j$ is observed firing. The effective
sample size for that edge is approximately the number of bins in which $j$ is
active,
\begin{equation}
  \neff(j) \;\approx\; f_j \,\Trec,
\end{equation}
with $f_j$ the firing rate of neuron $j$ and $\Trec$ the record duration. At
the median rate $f_j \approx 2.5$\,Hz this gives $\neff \approx 1{,}500$
presynaptic spikes at $10$ minutes and $\approx 9{,}000$ at $60$ minutes; for
the slowest neurons ($1$\,Hz) the $10$-minute count drops to $\approx 600$.
The slowest presynaptic neurons therefore bound the worst-case edge recovery.

For maximum-likelihood estimation of GLM coefficients, the standard error of
each weight scales as
\begin{equation}
  \mathrm{SE}(\Ahat_{ij}) \;\propto\; \frac{1}{\sqrt{\neff(j)}}.
\end{equation}
Reducing the record from $60$ to $10$ minutes is a $6\times$ reduction in data,
predicting
\begin{equation}
  \frac{\mathrm{SE}_{10}}{\mathrm{SE}_{60}}
  \;\approx\; \sqrt{6} \;\approx\; 2.45.
\end{equation}
This is the quantitative \emph{null hypothesis} for the duration sweep:
magnitude-estimation error should grow by $\approx\!2.45\times$ from $60$ to
$10$ minutes if sampling noise alone is responsible. Larger growth indicates an
additional failure mode (e.g.\ EM converging to a worse optimum, or support
errors corrupting the surviving magnitudes); smaller growth indicates a
regularization floor where $\lambda_3$ dominates sampling noise and duration
matters little. \textbf{Important caveat:} support recovery does \emph{not}
follow this smooth law. Weak edges sit near a detection threshold; as $\neff$
falls, edges near that threshold drop out abruptly, so recall (especially of
weak edges) degrades in a step-like fashion while precision can collapse from
proliferating noise-floor false positives.

% ===============================================================
\section{Part II: Comparing Support}
\label{sec:support}
% ===============================================================

The predicted support is the off-diagonal binary mask
$\hat E_{ij} = \mathbf{1}\{|\Ahat_{ij}| > \minW\}$; ground truth is $\Etrue$.
On the off-diagonal entries define the confusion counts
\begin{equation}
  \mathrm{TP},\;\mathrm{FP},\;\mathrm{FN},\;\mathrm{TN}
\end{equation}
(edge present in both / only in $\hat E$ / only in $\Etrue$ / absent in both).
Because the network is sparse ($\sim 5$--$20\%$ of entries are edges), the
overwhelming majority of entries are TN, and any metric that rewards TN will
read optimistically. We therefore favor edge-focused metrics and always report
the raw counts $(\mathrm{TP},\mathrm{FP},\mathrm{FN})$ so that a change in any
summary can be attributed to missed edges (FN) versus hallucinated edges (FP),
which are distinct failure modes.

\subsection{Candidate Scalar Metrics}
All are evaluated at the fixed $\minW = 0.02$.

\paragraph{Jaccard index (primary).}
\begin{equation}
  J = \frac{\mathrm{TP}}{\mathrm{TP}+\mathrm{FP}+\mathrm{FN}}.
\end{equation}
Ignores TN entirely; answers ``of all edges either matrix claims, what
fraction do both agree on.'' Most conservative and most edge-focused of the
candidates.

\paragraph{$F_1$ score.}
\begin{equation}
  F_1 = \frac{2\,\mathrm{TP}}{2\,\mathrm{TP}+\mathrm{FP}+\mathrm{FN}}.
\end{equation}
Also ignores TN. $F_1$ and $J$ are monotone transforms of one another,
$J = F_1/(2-F_1)$, so they rank fits identically; $J$ is simply harsher on the
same scale. Retained for comparison only.

\paragraph{Precision and Recall.}
\begin{equation}
  \mathrm{Prec}=\frac{\mathrm{TP}}{\mathrm{TP}+\mathrm{FP}},\qquad
  \mathrm{Rec}=\frac{\mathrm{TP}}{\mathrm{TP}+\mathrm{FN}}.
\end{equation}
Reported as the decomposition underlying $J$ and $F_1$: recall tracks loss of
real edges as $\neff$ falls, precision tracks the proliferation of noise-floor
false positives.

\paragraph{Matthews correlation coefficient (secondary).}
\begin{equation}
  \mathrm{MCC} =
  \frac{\mathrm{TP}\cdot\mathrm{TN} - \mathrm{FP}\cdot\mathrm{FN}}
       {\sqrt{(\mathrm{TP}+\mathrm{FP})(\mathrm{TP}+\mathrm{FN})
              (\mathrm{TN}+\mathrm{FP})(\mathrm{TN}+\mathrm{FN})}}.
\end{equation}
Uses all four cells, including TN, so it reads higher than $J$ in the sparse
regime; retained for comparability with the broader literature. Reports the
overall correlation between the two binary matrices.

We will run all of the above initially and discard the least discriminating
once empirical sensitivity is known.

\subsection{Block-Resolved Support (E/I structure)}
\label{sec:blocks}
The neuron ordering induces four meaningful blocks, indexed by presynaptic type
(\emph{column} block, by the from-$j$ convention) and postsynaptic type
(\emph{row} block):
\begin{equation}
  \text{E}\!\to\!\text{E},\quad
  \text{E}\!\to\!\text{I},\quad
  \text{I}\!\to\!\text{E},\quad
  \text{I}\!\to\!\text{I}.
\end{equation}
Each scalar metric of Section~\ref{sec:support} is recomputed within each block.
Because inhibitory neurons are a minority ($\sim 20\%$) and typically carry
different effective drive, we expect $\text{I}\!\to\!\cdot$ edges to recover
more poorly; the block table makes this circuit-meaningful rather than
arbitrary.

\subsection{Spatially Coarse-Grained Support}
\label{sec:grid}
Using the known $(x,y)$ coordinates, overlay a coarse $G\times G$ grid on the
neuron positions. Let $g(i)\in\{1,\dots,G^2\}$ be the grid cell of neuron $i$,
and partition neurons additionally by type $\tau(i)\in\{\mathrm{E,I}\}$. Define
a \emph{lumped} (super-node) connectivity between ordered cell-type pairs by
counting only edges that leave one cell and enter another. For source cell $a$
of type $\sigma$ and target cell $b$ of type (any, or fixed), the lumped true
and recovered edge counts are
\begin{align}
  \mathcal{N}^{\mathrm{true}}_{(a,\sigma)\to b}
  &= \sum_{\substack{j:\,g(j)=a,\,\tau(j)=\sigma \\ i:\,g(i)=b,\; i\neq j}}
     (\Etrue)_{ij}, \\
  \mathcal{N}^{\mathrm{hat}}_{(a,\sigma)\to b}
  &= \sum_{\substack{j:\,g(j)=a,\,\tau(j)=\sigma \\ i:\,g(i)=b,\; i\neq j}}
     \mathbf{1}\{|\Ahat_{ij}|>\minW\}.
\end{align}
Restricting the source sum to a single type $\sigma$ keeps excitatory and
inhibitory super-edges separate, as intended. Comparison at the super-node
level can use either a count-agreement measure,
\begin{equation}
  \mathrm{err}_{(a,\sigma)\to b}
  = \frac{\bigl|\mathcal{N}^{\mathrm{hat}}_{(a,\sigma)\to b}
                - \mathcal{N}^{\mathrm{true}}_{(a,\sigma)\to b}\bigr|}
         {\mathcal{N}^{\mathrm{true}}_{(a,\sigma)\to b}+1},
\end{equation}
or a fine-grained Jaccard restricted to the edge set between the two cells.
Rendering $\mathrm{err}_{(a,\sigma)\to b}$ as a spatial heatmap reveals whether
recovery quality varies across the tissue. \textbf{Interpretive caution:} in
the current generator the binary support is spatially i.i.d., so $\Etrue$ has no
intrinsic spatial structure. Any spatial pattern in recovery is therefore an
\emph{inference artifact}, and is expected to be explained almost entirely by
the spatial distribution of presynaptic firing rates rather than by position
per se (Section~\ref{sec:strat}). The grid view is thus a descriptive lens; the
mechanistic axis is firing rate.

\subsection{Stratified Recovery}
\label{sec:strat}
The mechanistic drivers of edge recoverability are presynaptic firing rate and
true edge magnitude. We therefore stratify recall:
\begin{itemize}[nosep]
  \item by presynaptic rate $f_j$ (binned),
  \item by true magnitude $|Atrue_{ij}|$ (binned),
  \item jointly, as a two-dimensional recall heatmap over
        $(f_j,,|Atrue_{ij}|)$.
\end{itemize}
The joint heatmap is the single most explanatory support diagnostic: it shows
exactly which edges become unrecoverable as duration shrinks, and it should
account for any apparent spatial or block-wise pattern.

% ===============================================================
\section{Part III: Comparing Magnitude}
\label{sec:magnitude}
% ===============================================================

Magnitude comparison is restricted to the true support
($Etrue_{ij}=1$, $i\neq j$) and uses the unit-$\rho$ matrices
$Ahat^{\star},Atrue^{\star}$ of Section~\ref{sec:order}. False-positive edges
are the concern of the support metrics and are excluded here, except where
noted (the global-vs-support Frobenius gap below). Because the true edges span
both signs (positive E, negative I), a metric computed across all signs is
partly rewarded for merely separating E from I---the easy part. To isolate
magnitude fidelity from sign, the primary forms of each metric are computed
\emph{within} the same-sign edge sets (within-E and within-I), with the global
(both-sign) form reported as a secondary, sign-inflated reference.

\subsection{Rank Fidelity: Spearman Correlation}
Over the chosen edge set, the Spearman correlation $\rho_S$ between
$\{Atrue^{\star}_{ij}\}$ and $\{Ahat^{\star}_{ij}\}$ measures preservation of
the \emph{ordering} of edge strengths:
\begin{equation}
  \rho_S = \mathrm{corr}\bigl(\mathrm{rank}(Atrue^{\star}),\;
                              \mathrm{rank}(Ahat^{\star})\bigr).
\end{equation}
Computed within-E, within-I, and globally. Spearman is invariant to any
monotone rescaling, so it is unaffected by the unit-$\rho$ normalization---a
useful internal consistency check (a rescaling bug moves Pearson but leaves
Spearman fixed). It weights all ranks equally and is therefore sensitive to the
fidelity of \emph{weak} edges.

\subsection{Magnitude Error: Relative Frobenius Norm}
Two complementary versions are reported.
\begin{align}
  \mathrm{RFE}_{\mathrm{supp}}
  &= \frac{\bigl\|(Ahat^{\star}-Atrue^{\star})\odot \Etrue\bigr\|_F}
          {\bigl\|Atrue^{\star}\odot \Etrue\bigr\|_F}, \\[4pt]
  \mathrm{RFE}_{\mathrm{glob}}
  &= \frac{\bigl\|Ahat^{\star}-Atrue^{\star}\bigr\|_F}
          {\bigl\|Atrue^{\star}\bigr\|_F},
\end{align}
both over off-diagonal entries. $\mathrm{RFE}_{\mathrm{supp}}$ measures
magnitude error on real edges; $\mathrm{RFE}_{\mathrm{glob}}$ additionally
includes the mass placed on false-positive entries. The gap
$\mathrm{RFE}_{\mathrm{glob}}-\mathrm{RFE}_{\mathrm{supp}}$ is a single-number
summary of spurious recovered mass and is expected to grow as duration shrinks.
The Frobenius norm is dominated by the largest-magnitude edges and is thus
sensitive to the \emph{strong} edges---complementary to Spearman.

\subsection{Systematic Shrinkage: Regression Slope}
Fit, over the chosen edge set, the linear relation
\begin{equation}
  Ahat^{\star}_{ij} = \beta\,Atrue^{\star}_{ij} + \gamma,
\end{equation}
reporting the slope $\beta$ and intercept $\gamma$ within-E and within-I, plus a
zero-intercept variant ($\gamma\equiv 0$). The slope $\beta$ is the shrinkage
diagnostic. After unit-$\rho$ rescaling the spectral-constraint mismatch is
removed, so a residual $\beta\neq 1$ is attributable to $\ell_1$ shrinkage and
finite-sample estimation bias; $\ell_1$ is expected to push $\beta<1$, more so
for the noisier inhibitory edges. A nonzero $\gamma$ on the support-restricted
fit flags a magnitude offset, which is physically unexpected (no edge implies
no systematic shift).

\subsection{The Underlying Scatter}
The three scalars above are summaries of one object: the scatter of
$Ahat^{\star}_{ij}$ against $Atrue^{\star}_{ij}$ over the true support,
colored by presynaptic rate $f_j$. This plot reveals shrinkage (slope),
heteroscedasticity (weak/slow edges noisier), and outliers (slow-presynaptic
edges) at a glance, and should be inspected before interpreting the scalars.

\subsection{Sensitivity Complementarity (Summary)}
The three magnitude metrics are deliberately complementary in what they
emphasize:
\begin{center}
\begin{tabular}{lll}
\toprule
Metric & Emphasizes & Scale dependence \\
\midrule
Spearman $\rho_S$ & ordering of all edges (incl.\ weak) & invariant \\
Relative Frobenius & large-magnitude (strong) edges & requires unit-$\rho$ \\
Regression slope $\beta$ & systematic shrinkage bias & requires unit-$\rho$ \\
\bottomrule
\end{tabular}
\end{center}

% ===============================================================
\section{Headline Deliverable}
\label{sec:deliverable}
% ===============================================================

The duration study is summarized by curves of the retained primary metrics
\begin{equation}
  \{\,J,\; \rho_S^{\,\mathrm{E}},\; \rho_S^{\,\mathrm{I}},\;
  \mathrm{RFE}_{\mathrm{supp}},\; \beta^{\,\mathrm{E}},\; \beta^{\,\mathrm{I}}\,\}
  \quad\text{vs.}\quad D\in\{10,\dots,60,120\}\text{ min},
\end{equation}
each normalized to its $120$-minute value and averaged over seeds with error
bars. The magnitude-error curve is overlaid with the
$\sqrt{6}\approx 2.45$ sampling-noise reference (Section~\ref{sec:neff});
deviations from this reference, together with the step-like behavior of recall
in the stratified heatmap (Section~\ref{sec:strat}), constitute the principal
findings on information loss between the $10$- and $60$-minute records.

\end{document}
